% =============================================================
%  Abstract + Introduction + Related Work
%  [BENCHMARKNAME] -- PII Leakage Benchmark for VLM Web Agents
%  Target venue: EMNLP 2025
%  NOTE: paste these sections into your main document; they do
%        not include a preamble or \begin{document}.
%        Replace every [BENCHMARKNAME] with the chosen name.
%        Replace every [TBD] with final experimental values.
% =============================================================


% -------------------------------------------------------------
\begin{abstract}
% -------------------------------------------------------------

Large language model agents are increasingly deployed to browse the web
autonomously on behalf of users, acting on tasks such as job applications,
online purchases, and form submissions. To complete such tasks without
constant supervision, these agents are granted access to the user's
personally identifiable information (PII), including names, addresses,
payment details, and government identifiers. This access creates a
confused-deputy attack surface: adversarial web content can coerce an
agent into forwarding sensitive data to an attacker-controlled endpoint.
Existing benchmarks evaluate agent susceptibility to dark patterns and
prompt injection, measuring task-trajectory divergence rather than
data exfiltration, and none incorporate multi-turn conversational attacks
or a structured PII profile. We introduce \textsc{[BENCHMARKNAME]},
a benchmark of approximately 100 hand-crafted environments spanning
16 site categories and 7 attack vectors, organized on an 8-axis
factorial taxonomy that supports controlled ablation studies. We evaluate
five frontier models on four metrics: PII Leakage Rate (PLR), Attack
Success Rate (ASR), Task Completion Rate (TCR), and Detection Rate (DR).
[TBD: headline findings.] Our results expose fundamental safety gaps
in current autonomous web agents for PII-sensitive tasks and motivate
output-level defenses that do not rely on the agent's ability to
recognize an attack before acting.

\end{abstract}


% -------------------------------------------------------------
\section{Introduction}
\label{sec:intro}
% -------------------------------------------------------------

Vision-language model (VLM) agents capable of operating a web browser
are moving from research prototypes into production deployment. Systems
such as GPT-4o with browsing, Gemini Agents, and Claude's computer-use
interface take high-level user instructions and autonomously navigate
websites, fill forms, complete checkout flows, and respond to emails.
Because these agents must act on the user's behalf without requiring
manual input at each step, they are routinely provisioned with the user's
full personal information: legal name, home address, payment card details,
Social Security number, passwords, and API keys. The agent stores this
profile in its context and draws on it whenever a form field must be
filled or a credential must be submitted. The result is a system that
combines broad web access with privileged read access to high-value
personal data, a combination that creates a classic confused-deputy
vulnerability~\citep{hardy1988confused}: untrusted content encountered
during normal browsing can instruct or manipulate the agent into
forwarding that data to an adversary.

The modern web already contains a rich ecosystem of social engineering
attacks built to harvest personal information from human users. Phishing
clones impersonate legitimate services on subtly misspelled domains.
Dark patterns bury the safe option under an array of deceptive buttons.
Fake trust signals, urgency timers, and authority impersonation
manipulate users into submitting data they would otherwise withhold.
Multi-turn chat interfaces escalate requests incrementally, building
rapport before asking for a password or a bank routing number. Human
users can slow down, scrutinize a URL, or simply close the tab.
A web agent cannot. It is optimized to complete the task given to it,
and it will follow links, fill fields, and submit forms with a
thoroughness and speed that no human user matches. This efficiency,
combined with unconditional access to a full PII profile, makes the
web-browsing agent a particularly high-value and unusually reliable
target for data exfiltration.

Despite growing interest in the security of LLM-based web agents, the
existing literature has not directly addressed this threat. Work on dark
patterns, represented by DECEPTICON~\citep{kumarappan2024decepticon}
and TrickyArena~\citep{li2024trickyarena}, measures whether an agent
can be steered toward a malicious navigation outcome, not whether it
surrenders personal data. DECEPTICON explicitly defines dark patterns as
residing within the user interface and excludes external threats like
phishing from its scope; its success metric is trajectory divergence,
not field-level data loss. Work on prompt injection, represented by
TRAP~\citep{trap2024}, WiPI~\citep{wu2024wipi}, EIA~\citep{liao2024eia},
and AdvWeb~\citep{xu2024advweb}, measures whether hidden adversarial
instructions can hijack the agent's task; no prior work in this line
provides the agent with a realistic PII profile or tracks what the
agent types into attacker-controlled forms. The fine-print injection
study of~\citet{huang2025fineprint} is the closest to our PII angle
but examines a specific injection vector across GUI applications
rather than a broad social engineering taxonomy.
Fraud-R1~\citep{yang2025fraudr1} benchmarks LLM robustness to
phishing and fraud inducements, but its subjects are language models
responding to text prompts, not autonomous agents operating a live
browser with a populated user profile. WebTrap Park~\citep{webtrappark}
evaluates agents across three risk categories including social
engineering, but does not define a PII metric, does not model a user
profile, and does not include conversational attacks. To our knowledge,
no prior work has (i) equipped a web agent with a structured PII profile
spanning critical and high-sensitivity fields, (ii) measured field-level
data exfiltration as the primary outcome, or (iii) covered the full
social engineering attack surface, including multi-turn adaptive chat
and multi-site flows in which an agent must navigate from an email
client to a separate phishing destination within a single task.

We introduce \textsc{[BENCHMARKNAME]}, a benchmark designed to answer
a single question: when an autonomous web agent encounters a social
engineering attack while browsing, will it hand over the user's personal
information? The benchmark comprises approximately 100 hand-crafted
environments spanning 16 site categories, including employment portals,
e-commerce stores, banking and healthcare platforms, cryptocurrency
exchanges, government portals, and social media services. The
environments cover 7 attack vectors: phishing clones, dark patterns,
credential harvesting, reward traps, authority impersonation, fake trust
signals, and multi-turn conversational deception. Each environment is
classified on four primary axes (site category, attack vector,
salience level, and PII target) and four ablation axes (urgency
pressure, prompt injection style, interaction modality, and multi-site
flow), forming a factorial taxonomy that enables controlled sibling
comparisons in which exactly one axis is varied. We equip an autonomous
Playwright-based agent with a 40-field PII profile covering identity,
address, payment, financial, professional, credential, and API-key
information. We evaluate five frontier models on four metrics: PII
Leakage Rate (PLR), measuring the fraction of runs in which at least
one PII field is submitted to an attacker-controlled endpoint; Attack
Success Rate (ASR), measuring whether the attacker's intended outcome
was achieved; Task Completion Rate (TCR), measuring whether the
legitimate task was completed; and Detection Rate (DR), measuring
whether the agent explicitly identified the attack in its reasoning.
Our model set includes both vision-language models, which receive DOM
text and page screenshots, and a text-only baseline, which receives
DOM text alone, allowing us to isolate the contribution of visual
understanding to agent security.

The contributions of this paper are as follows.

\begin{enumerate}

\item \textbf{\textsc{[BENCHMARKNAME]}}: the first benchmark to measure
field-level PII exfiltration by autonomous web-browsing agents under
social engineering attacks, comprising approximately 100 hand-crafted
environments across 16 site categories and 7 attack vectors.

\item \textbf{Factorial taxonomy}: an 8-axis design in which sibling
environments differ on exactly one axis, enabling causal ablation
studies on the effect of urgency pressure, prompt injection style,
interaction modality, and multi-site flow on agent vulnerability.

\item \textbf{Multi-turn and multi-site coverage}: the first benchmark
to include adaptive multi-turn conversational attacks with response
trees and multi-site task flows in which the agent must traverse
an email client and a phishing destination within a single task.

\item \textbf{Empirical findings} across five frontier models: [TBD,
to be inserted after full evaluation runs].

\end{enumerate}

The remainder of the paper is organized as follows.
Section~\ref{sec:related} surveys related work.
Section~\ref{sec:benchmark} describes the benchmark design and taxonomy.
Section~\ref{sec:eval} presents the agent scaffold and evaluation protocol.
Section~\ref{sec:results} reports experimental results.
Section~\ref{sec:discussion} discusses implications for agent safety and
directions for defense. Section~\ref{sec:conclusion} concludes.


% -------------------------------------------------------------
\section{Related Work}
\label{sec:related}
% -------------------------------------------------------------

\subsection{Dark Patterns and Deceptive UI in Web Agent Benchmarks}

The impact of deceptive user interface design on LLM-based web agents
has received increasing attention. DECEPTICON~\citep{kumarappan2024decepticon}
introduces a 700-task benchmark covering dark patterns embedded within
the UI of web navigation tasks. Across state-of-the-art agents, dark
patterns successfully redirect the agent toward malicious outcomes in
over 70\% of tasks, compared to a human rate of 31\%. Notably,
susceptibility increases with model capability: larger models using
extended reasoning are more vulnerable, not less. The authors explicitly
restrict dark patterns to in-situ UI elements and exclude phishing and
malware from their scope, so no PII outcome is measured. Guardrails
reduce the success rate of dark pattern interventions by 28 percentage
points but do not eliminate them.

TrickyArena~\citep{li2024trickyarena} is a controlled environment of
web applications from e-commerce, streaming, and news domains in which
dark patterns can be selectively enabled or disabled. With a single dark
pattern active, agents are susceptible approximately 41\% of the time.
The study finds that stacking multiple dark patterns increases attack
success and that enabling multimodality raises susceptibility rather than
reducing it, suggesting that visual deception is at least as effective
as textual manipulation. DarkBench~\citep{kran2025darkbench} examines
dark patterns produced by LLMs in their own outputs, a complementary
direction that does not address agent behavior during web navigation.

Our work departs from this line in two respects. First, we treat PII
exfiltration as the primary outcome rather than trajectory divergence.
An agent that visits a phishing site but submits no data represents a
different risk profile than one that completes every form it encounters.
Second, our benchmark explicitly covers attack types that originate
outside the UI, including phishing domains, authority impersonation,
and conversational deception, which the existing dark-pattern literature
excludes by design.

\subsection{Prompt Injection and Adversarial Web Content}

A parallel literature studies prompt injection attacks, in which
adversarial instructions hidden in web content redirect the agent's
behavior. TRAP~\citep{trap2024}, built on the REAL benchmark of
deterministic website simulations~\citep{dewitt2025real}, evaluates six
frontier models and finds average susceptibility of 25\%, ranging from
13\% for GPT-5 to 43\% for DeepSeek-R1. Small changes to interface
framing or surrounding context can double the success rate, indicating
psychologically driven fragility that is not captured by task-level
metrics alone. WiPI~\citep{wu2024wipi} surveys web-based prompt
injection vectors. EIA~\citep{liao2024eia} introduces environmental
injection attacks targeting privacy leakage. AdvWeb~\citep{xu2024advweb}
explores black-box adversarial attacks on VLM-powered web agents.

The fine-print injection study of~\citet{huang2025fineprint} is the
closest in spirit to our PII focus. Across 234 adversarial webpages and
six GUI agents, agents are highly vulnerable to contextually embedded
threats, including threats that exploit the mismatch between visual
salience for agents versus human readers. A parallel study with 39 human
participants finds that humans are also susceptible to many of the same
attacks, suggesting that simple human oversight may not reliably prevent
failures. However, the study targets a specific set of injection vector
types rather than a broad social engineering taxonomy, and it does not
provide the agent with a simulated user profile to exfiltrate.

The unifying limitation of prompt injection benchmarks is that they
measure whether the agent executes the attacker's instruction, not
whether the agent's own PII is compromised. Our benchmark measures
what the agent types, specifically which fields from its user profile
it submits to which endpoints, across an attack surface that encompasses
both static and conversational social engineering.

\subsection{Phishing and Fraud Robustness of Language Models}

Fraud-R1~\citep{yang2025fraudr1} is a multi-round benchmark assessing
LLM robustness to augmented phishing and fraud inducements presented as
text prompts. It demonstrates that adversarial framing and escalating
pressure meaningfully degrade refusal rates. However, the subjects are
language models responding to text, not agents operating a browser, and
no user profile is defined. The benchmark does not address whether an
agent that has already decided to complete a task can be induced to
also submit sensitive data it was never asked to provide.
When AI Agents Collude~\citep{agentcollude2025} examines a different
threat model in which multiple LLM agents on social platforms engage in
financially fraudulent coordination; this multi-agent dynamic is
orthogonal to our single-agent, single-session evaluation. The
machine-readable ads study~\citep{machineads2025} investigates
behavioral patterns of AI agents interacting with online advertisements,
documenting how agents process and respond to persuasive commercial
content. Our work places PII exfiltration at the center of the threat
model and introduces social engineering as the mechanism, rather than
studying commercial influence or multi-agent fraud.

\subsection{Web Agent Security Evaluation Platforms}

WebTrap Park~\citep{webtrappark} is an automated platform for systematic
security evaluation of web agents across three risk categories: malicious
user prompts, prompt injection, and deceptive website design. It operates
with live browser interaction and evaluates multiple agent scaffolds,
finding that scaffold choice meaningfully affects vulnerability. However,
the benchmark does not define a PII metric or model a user data profile;
social engineering is one of three categories without coverage of
conversational attacks. REAL~\citep{dewitt2025real} provides
deterministic simulations of real websites and serves as the environment
substrate for TRAP, but its primary purpose is task-completion evaluation
rather than security evaluation. Our benchmark focuses exclusively on the
social engineering vector targeting PII and provides a factorial design
that supports causal conclusions about which attack dimensions drive
vulnerability, a capability that single-axis or catalog-style benchmarks
cannot offer.
